% AICE 分散ベンチマーク報告 — N-Queens と 食事する哲学者
% コンパイル: lualatex aice-bench-report.tex
\documentclass[a4paper,11pt]{ltjsarticle}
\usepackage{luatexja}
\usepackage[margin=24mm]{geometry}
\usepackage{booktabs}
\usepackage{listings}
\usepackage{xcolor}
\usepackage{enumitem}
\usepackage[hidelinks]{hyperref}

\definecolor{codebg}{gray}{0.96}
\lstset{basicstyle=\ttfamily\small,backgroundcolor=\color{codebg},frame=single,
  rulecolor=\color{gray!40},breaklines=true,columns=fullflexible,keepspaces=true,
  xleftmargin=4pt,xrightmargin=4pt,aboveskip=6pt,belowskip=6pt}
\newcommand{\cmd}[1]{\texttt{#1}}

\title{\textbf{AICE 分散ベンチマーク報告\\ N-Queens と 食事する哲学者\\ (Mac + Embedded Xinu / Pi 5)}}
\author{AICE プロジェクト}
\date{2026年6月8日}

\begin{document}
\maketitle

\begin{abstract}
\noindent
本日利用可能な計算資源（管理ノード = 本 Mac 1 台、ワーカーノード = Embedded
Xinu を載せた Raspberry Pi 5 1 台）から成る AICE システム上で、型推論クリーンな
AIPL で記述した2つの古典的問題——\textbf{N-Queens}（探索）と\textbf{食事する哲学者}
（並行）——を分散実行し、ベンチマークを取得した。N-Queens は第1行の列で探索空間を
分割して Mac と Pi 5 に振り分け、哲学者は5人×50食（計250食）に要する時間を測定した。
いずれもガベージコレクション(GC)を稼働させた状態で計測している。
\end{abstract}

\tableofcontents

%======================================================================
\section{実験環境}
\begin{table}[h]\centering
\begin{tabular}{ll}
\toprule
役割 & 実体 \\
\midrule
管理ノード & 本 Mac（AIPL OCaml/Python ランタイム、aipl2c コンパイラ）\\
ワーカーノード & Raspberry Pi 5 / BCM2712、Embedded Xinu（ベアメタル）\\
Xinu IP & \texttt{192.168.3.101}（HTTP）\\
Xinu 実行系 & on-device C-JIT（\texttt{/cc} 経由、AArch64 にコンパイルして実行）\\
GC & Xinu: 約8秒周期の自動アクタGC + \texttt{/api/actors-gc}。Mac: ランタイム内蔵GC \\
\bottomrule
\end{tabular}
\caption{AICE 構成（本日）}
\end{table}

\noindent
役割分担は「管理ノード(Mac)が AIPL ロジックの実行とオーケストレーション、Pi 5 が
計算ワーカー」である。AIPL は型推論クリーン（\cmd{aipl2c --check} が型エラーなしで通過）
であり、Pi 5 へは \cmd{aipl2c} で C に変換して \cmd{/cc} に送る。

%======================================================================
\section{N-Queens（探索問題）}

\subsection{記述と分散方式}
N-Queens は $N\times N$ 盤に互いに利かない $N$ 個のクイーンを置く配置数を数える問題。
AIPL で再帰的バックトラック（\cmd{safe}/\cmd{nq\_count}）として記述し、
\textbf{第1行のクイーンの列 $k$ で探索空間を $N$ 個の部分問題に分割}した。各 $k$ は独立に
数えられるので、Mac と Pi 5 へ自由に振り分けられる。

\begin{lstlisting}[language=C]
function nq_count(cols, n, row) {     // AIPL（型推論クリーン）
  if (row == n) { return 1; }
  var total = 0; var c = 0;
  while (c < n) do {
    if (safe(cols, row, c)) { cols[row] = c;
      total = total + nq_count(cols, n, row + 1); }
    c = c + 1; }
  return total; }
\end{lstlisting}

\paragraph{$N$ の選択.}Pi 5 の C-JIT はループ毎に挿入される実行時間ガード
(\cmd{cc\_tick}, 約 250\,ms の暴走打ち切り)を持つため、1 回の \cmd{/cc} 呼び出しで
完了する計算量に上限がある。実測の結果、\textbf{部分問題を1列ずつ（$k$ 単位で）
個別に呼ぶ}ことで $N{=}9$（厳密解 352）まで正しく完了することを確認した
（$N{=}10$ は一部の部分問題が打ち切られ不正解になる）。よって本ベンチは $N{=}9$ とする。

\subsection{結果（$N=9$、厳密解 352）}
\begin{table}[h]\centering
\begin{tabular}{lrrrr}
\toprule
構成 & 解の数 & 合計(s) & Mac(s) & Pi5(s) \\
\midrule
Mac 単独（AIPL Py·I）        & 352 & 1.70 & 1.70 & 0.00 \\
Pi 5 単独（cc JIT）          & 352 & 1.12 & 0.00 & 1.12 \\
Mac=0 + Pi5=9               & 352 & 1.17 & 0.00 & 1.17 \\
Mac=2 + Pi5=7               & 352 & 0.92 & 0.57 & 0.91 \\
\textbf{Mac=4 + Pi5=5}      & \textbf{352} & \textbf{0.91} & 0.91 & 0.65 \\
Mac=6 + Pi5=3               & 352 & 1.21 & 1.21 & 0.39 \\
Mac=9 + Pi5=0               & 352 & 1.71 & 1.71 & 0.00 \\
\bottomrule
\end{tabular}
\caption{N-Queens $N{=}9$ の負荷分散（第1行の列を Mac と Pi 5 に分配、並列実行）。
全構成で解の数 352 を正しく得た。}
\end{table}

\subsection{考察}
\begin{itemize}[noitemsep]
  \item 単独では Pi 5（1.12\,s）が Mac の AIPL インタプリタ実行（1.70\,s）より速い。
        Pi 5 は AArch64 ネイティブにコンパイルして実行するため。
  \item \textbf{分散により両者単独を上回る高速化}が得られ、最良は Mac=4+Pi5=5 の
        \textbf{0.91\,s}（Mac 単独比 1.87 倍、Pi5 単独比 1.23 倍）。両ノードが
        並列に部分問題を消化するため。
  \item 偏った分割（Mac=6+Pi5=3 等）は遅い側（Mac）が律速となり改善が鈍る。
        最適点は各ノードの単位速度に応じた配分付近にある。
\end{itemize}

%======================================================================
\section{食事する哲学者（並行問題）}

\subsection{記述と測定}
5 人の哲学者が円卓で隣り合う2本のフォークを取って食事する古典的並行問題。
AIPL では哲学者とフォークをアクタとして記述し、\cmd{sender} ベースのフォーク獲得で
デッドロックを回避する。\textbf{各哲学者が 50 食、計 250 食}を完了するまでの時間を測定した。
Pi 5 側はアクタ群のクロストーク負荷が大きいため、同一アルゴリズムの逐次 C シミュレータ
（毎ラウンド、両フォークが空く哲学者が1食ずつ進む）を \cmd{/cc} で実行した。

\subsection{結果（5人 × 50食 = 250食）}
\begin{table}[h]\centering
\begin{tabular}{lrr}
\toprule
構成 & 完了食数 & 時間(ms, 中央値) \\
\midrule
Mac 単独（AIPL アクタ, Py·I） & 250 & 41（41/41/38）\\
Pi 5 単独（cc JIT 逐次C）      & 250 & 69（61/71/69）\\
\bottomrule
\end{tabular}
\caption{食事する哲学者、5人×50食。両者とも 250 食を完了。Pi 5 の時間は Mac 側で
計測した HTTP 往復込みの wall time で、その大半は cc JIT のコンパイル時間
（実際の 250 食前進のネイティブ実行は数十マイクロ秒）。}
\end{table}

\subsection{考察}
\begin{itemize}[noitemsep]
  \item Mac の AIPL アクタスケジューラは HTTP もスリープも介さないため約 40\,ms と速い。
  \item Pi 5 単独は約 69\,ms。内訳の支配項は \cmd{/cc} の JIT コンパイル時間であり、
        実行そのものは極小。すなわち Pi 5 では「投入のコスト」が「計算のコスト」を
        上回る軽量タスクであることを示す。
  \item この特性から、哲学者のような\textbf{細粒度・低計算量タスクは投入オーバーヘッドが
        支配的}で分散の利得が出にくい。一方 N-Queens のような\textbf{粗粒度・高計算量タスクは
        分散で素直に高速化}する（第3節）。AICE における配置ポリシは、このタスク粒度を
        考慮して Mac/Xinu を選ぶべきである。
\end{itemize}

%======================================================================
\section{ガベージコレクション(GC)}
両ベンチとも GC を稼働させた。
\begin{itemize}[noitemsep]
  \item \textbf{Pi 5}：約 8 秒周期で自動アクタGC（20 秒以上アイドルのアクタを回収）が常時稼働。
        加えてベンチ前後に \cmd{/api/actors-gc} を明示実行した。N-Queens は plain C で
        アクタを生成しないため回収対象 0（\texttt{no resident actors}）。哲学者の Pi 5 側も
        逐次 C のため同様。
  \item \textbf{Mac}：AIPL ランタイム内蔵の GC が常時稼働。哲学者の Mac 側は 5 哲学者
        + 5 フォーク + Tally のアクタを生成し、実行後に回収される。
\end{itemize}
GC を有効にしたまま上記の時間が得られており、GC が本ワークロードの性能を阻害しないことを確認した。

%======================================================================
\section{結論}
\begin{itemize}[noitemsep]
  \item 型推論クリーンな AIPL で記述した N-Queens・哲学者を、Mac + Pi 5 の AICE 上で
        分散実行し、いずれも正しい結果（N-Queens 352、哲学者 250 食）を得た。
  \item \textbf{N-Queens（粗粒度）は分散で高速化}（最良 0.91\,s、Mac 単独比 1.87 倍）。
  \item \textbf{哲学者（細粒度）は投入オーバーヘッドが支配的}で、軽量タスクの分散は
        利得が小さい——タスク粒度に応じた配置の重要性を実証。
  \item GC を稼働させたまま計測しても性能劣化は見られなかった。
\end{itemize}

\vfill
\begin{center}\small AICE / Embedded Xinu Raspberry Pi 5 — 2026年6月8日\end{center}
\end{document}
